\documentclass{zhvt-classic}

% 使用本地源流明体（GenRyuMin2 TW）按文件路径加载，不依赖系统字体名
% 版框（内框／界栏／书口侧栏／鱼尾）改用与版心印章同一印泥棕色：
%   印章图片 yshs.png 实色 = RGB(135,68,52) = HTML 874434
\usepackage{xcolor}
\usepackage{tikz}
% sealbrown 双注册：xcolor 供 TikZ/pgf，L3 供 cls 版框
\definecolor{sealbrown}{HTML}{874434}
\ExplSyntaxOn
\color_set:nnn { sealbrown } { HTML } { 874434 }
\ExplSyntaxOff

\zhvtset{
  grid color       = sealbrown,
  main font        = GenRyuMin2TW-SB,
  main font path   = C:/Users/Administrator/texmf/fonts/opentype/public/GenRyuMin2TW/,
jiazhu font      = GenRyuMin2TW-M,
  jiazhu font path = C:/Users/Administrator/texmf/fonts/opentype/public/GenRyuMin2TW/,
  % 页码框容纳 4 字：正常页「九十七」+前置空 1 格=4 行恰好；
  % 100+ 页「一百零一」共 5 行超出，末字「一/二/三」被框裁掉（2026-10-04 实测），故放宽到 5
  shukou page number chars = 5,
  % 书口书名框高从默认 6 字减到 4 字，把省出的 2 字高拨给 chapter 框，
  % 使 6 字书名+卷名放进 9 字高的 chapter 框后不再压线（2026-10-05）
  shukou title chars = 4,
}

\title{史记}[韩信卢绾列传]
\author{汉·司马迁}
\maker{闲叟漫记·重製}


% ===== 索隐述赞样式（2026-11-08 四次定稿 · 愚大人亲手改定参数）=====
% ⚠ 竖排前提：自定义字体家族必须带 RawFeature={vertical:+vert}；
%   simhei.ttf 加 +vhal 会整体加载失败静默回退，只写 +vert；simfang 同理。
\newCJKfontfamily{\shuzanfont}[Path=C:/Windows/Fonts/,RawFeature={vertical:+vert}]{simhei.ttf}
\newCJKfontfamily{\shuzanbody}[Path=C:/Windows/Fonts/,RawFeature={vertical:+vert}]{simfang.ttf}
% 题头：152x30pt 圆角 6pt 棕底白字；baseline (0,2) 使整框落入界行格子不压版框线；
% 白字节点 x 不补偿、y +8.2pt（像素级墨迹实测 dx/dy≈0，L/R 留白约 20pt）。
\newcommand{\shuzanhead}[1]{%
  \begin{tikzpicture}[baseline={(0,2)}]
    \fill[sealbrown, rounded corners=6pt]
      (-76pt,-15pt) rectangle (76pt,15pt);
    \node[text=white, inner sep=0pt, anchor=center] at (0pt,8.2pt)
      {\shuzanfont\fontsize{23}{23}\selectfont #1};
  \end{tikzpicture}%
}
% ★ 修正书口卷名位置：chapter 框起始位置改回 cls 默认值，避免全局影响。
%   只在每个卷首页通过 \juan 命令里的空命令微调该页位置（安全，不会跨页污染）。
\ExplSyntaxOn
\AtBeginDocument{
  \dim_set:Nn \l__zhvt_shukou_chapter_title_start_dim
    { \l__zhvt_shukou_title_height_dim + \l__zhvt_font_size_dim * 2 }
  \dim_set:Nn \l__zhvt_shukou_page_number_start_dim
    { \l__zhvt_shukou_chapter_title_start_dim + \l__zhvt_shukou_chapter_title_height_dim }
}
\ExplSyntaxOff

\begin{document}

% ★★ 版心（书口）排布覆写：把书名放到鱼尾之下（2026-10-04）
%
%   cls 默认版心自上而下：书名（title 静态框，在鱼尾之上）→ 鱼尾 → 篇名（chapter 动态框）→ 页码
%   本模板改为：          鱼尾（位置不动）→ 书名 → 篇名 → 页码分割线 → 页码
%
%   手法：把 cls 放在 title 静态框里的书名，改放进 chapter 动态框（鱼尾之下）；
%         title 静态框留空。书名取 \l__zhvt_title_str，即 \title{...} 的第一参数，
%         所以换书时只改 \title，不必在此重复写书名。
%
%   ⚠ 三条必须遵守的写法（2026-10-04 实测）：
%     ① 汉字必须用 \hbox{...} 包住。flowfram 会把框内容装进一个宽 14.4mm（= 一个 line_sep
%        = 书口栏宽）的纵向 minipage 并按该宽度折行，而整页靠 PDF /Rotate 90 显示、
%        框内文字本身不变向 —— 不包 \hbox 时「围炉夜话」4 字会被逐字折成 4 行，
%        在版面上就成了横排 4 字并横向溢出版框。
%     ② 竖排下**一个 \hbox 就是一整列**。书名与篇名要同处一列，必须放进同一个 \hbox；
%        分成两个 \hbox 时，后一个会被当成新的一列、甩到隔壁栏去
%        （实测「序」跑到书口栏外侧 14.4mm 的正文列上）。中间的 \quad 在框内是横向空白，
%        显示出来正好是书名与篇名之间的纵向间隔。
%     ③ ★ 内容盒要写成 **与框等宽的定宽盒** `\hbox to \l__zhvt_shukou_chapter_title_height_dim{…\hfil}`。
%        实测：chapter 框是把内容盒沿纵向**居中**摆放的，内容宽度一变，书名就跟着漂：
%        自然宽度的 \hbox（书名 38.4mm）落在 x=157.2mm，加篇名后（57.6mm）落在 x=152.4mm，
%        目录页（无篇名）与正文页的书名就差半格。做成等宽盒后，无论有没有篇名，
%        内容固定从框顶起排，实测两页均落在 x=142.8mm，不再漂。
\ExplSyntaxOn
% 把书口 chapter 框（书名+篇名）起始位置下移 1 字高，避免 4 字卷名顶到鱼尾
% ⚠ cls 的 \__zhvt_calculate_dimensions: 在 \AtBeginDocument 里执行，会重算该 dim；
%   必须把覆写挂到 begin document 之后，否则会被二次覆盖（2026-10-05 实测）。
% (旧方案已移到导言区，直接修改 chapter 框起始位置)
% 卷标题命令 #1=书口篇名(鱼尾) #2=正文卷名（含卷序，卷名内空格用 ~）
%
% ★★ 目录页两项修订（2026-10-05，愚大人指出）★★
%   ① 页码起点统一对齐：
%      原状 —— 目录项写成 \par\hspace{4em}卷名\hspace{4em}页码，页码起点 = 卷名实宽 + 4em。
%      卷名实宽随字数变：「卷一~梁惠王上」6 字 +~、「卷七~离娄上」5 字 +~，
%      差 1 个字位 → 卷七~卷十 4 条的页码整体比其余 10 条高 27.21pt（一个字），看着上下不齐。
%      改法 —— 卷名补全角空格（宽 = 1 字 = 27.215pt）补足到 **7 字符位**再写进目录，
%      14 条页码起点遂全部落在同一条竖线上（起点对齐；末字随页码字数自然参差）。
%      ⚠ 只对目录里的副本 \l_tmpc_tl 补，正文卷名仍用 \l_tmpa_tl 原样 —— 正文像素不受影响。
%   ② 卷名整体上移半字位：前导 \hspace 由 4em 改 3.5em。
%      实测原状目录块在版框内偏下（天头留白 139.6pt / 地脚留白 110.7pt），
%      上移 0.5em = 13.6pt 后上下基本均分，视觉居中。
\NewDocumentCommand { \juan } { mm } {
	\clearpage
	\tl_set:Nn \l__zhvt_chapter_title_tl { #1 }
	\tl_set:Nn \l_tmpa_tl { #2 }
	\tl_set_eq:NN \l_tmpc_tl \l_tmpa_tl
	\int_step_inline:nn { 7 - \str_count:n { #2 } }
	  { \tl_put_right:Nn \l_tmpc_tl { 　 } }
	\tl_set:Nx \l_tmpb_tl { \zhnum { page } }
	\tl_put_left:Nn \l_tmpb_tl { \hspace { 4em } }
	\tl_put_left:Nn \l_tmpb_tl { \l_tmpc_tl }
	\tl_put_left:Nn \l_tmpb_tl { \par \hspace { 3.5em } }
	\addtocontents { toc } { \l_tmpb_tl }
	\skip_horizontal:n { \l__zhvt_font_size_dim * 4 }
	\l_tmpa_tl
	\par
	\thispagestyle { empty }
}
\cs_set_protected:Npn \__zhvt_set_shukou_content: {
	% ★ 鱼尾上方留白（title 静态框）嵌入闲叟钤印：
	%   整页 /Rotate 90，故此处把图用 \rotatebox{90} 转正为竖章；
	%   尺寸卡在版心条幅内（书内约 7mm 宽 × 33.8mm 高，留余量不压鱼尾）。
	\setstaticcontents * { title } {
		\vbox_to_ht:nn { \l__zhvt_baseline_skip_dim } {
			\vskip 14.2pt
			\hbox_to_wd:nn { \l__zhvt_shukou_title_height_dim } {
				\hfil
				\rotatebox [ origin = c ] { 90 } { \includegraphics [ width = 7mm , height = 33.8mm ] { yshs.png } }
				\hfil
			}
		}
	}
% ★ 版心（书口）统一改法（2026-11-16 愚大人定）：18pt＋框内居中＋书口文字同版框色
%   ① 原 27.2pt 时 7 字长篇名（如「廉颇蔺相如列传」）末字溢出 chapter 框、压住页码；
%      改 \fontsize{18pt}{18pt}，前后 \hfil 令字串在框内居中（不随篇名字数漂移）。
%   ② \vbox{\vskip 8.2pt} 定量下移，抵消缩字号后栏宽方向 8.2pt 上偏（像素实测 y 中心 543.3 = 原版）。
%      ⚠ \raisebox / \vspace* 在 flowfram 框内无效；换字号须同步换补偿：16pt→10.1、19pt→7.4。
%   ③ chapter 与 page~number 均置 \color{sealbrown}（HTML 874434），与版框／界行同色；正文仍黑。
%   ⚠ 页码字位、钤印、鱼尾、界行、正文一律未动；五部合订本同此口径。
	\setdynamiccontents * { chapter }
	{ \color{sealbrown} \vbox { \vskip 8.2pt \hbox to \l__zhvt_shukou_chapter_title_height_dim
		{ \hfil \fontsize{18pt}{18pt}\selectfont \l__zhvt_title_str \quad \l__zhvt_chapter_title_tl \hfil } } }
	\setdynamiccontents * { page~number } { \color{sealbrown} \quad \zhnum { page } }
}
\__zhvt_set_shukou_content:
\ExplSyntaxOff

\maketitle

% 单篇，无多卷目录，不渲染目录页

\juan{韩信卢绾列传}{卷九十三~韩信卢绾列传第三十三}
\hspace*{2em}韩王信者。故韩襄王孽孙也。长王尺五寸。及项梁之立楚後怀王也。燕。齐。赵。魏皆已前王。唯韩成有後。故立韩诸公子横阳君成为韩王。欲以抚定韩故地。项梁败死定陶。成奔怀王。沛公将兵击阳城。使张良以韩司徒降下韩故地。得信。以为韩将。将其兵从沛公入武关。\\
\hspace*{2em}沛公立为汉王。韩信从入汉中。乃说汉王曰。﹁项王王诸将近地。而王独远居此。此左迁也。士卒皆山东人。跂而望归。及其锋东乡。可以争天下。﹂汉王还定三秦。乃许信为韩王。先拜信为韩太尉。将兵略韩地。\\
\hspace*{2em}项籍之封诸王皆就国。韩王成以不从成功。不遣就国。更以为列侯。及闻汉遣韩信略韩地。乃令故项籍游吴时吴令郑上为韩王以距汉。汉二年。韩信略定韩十馀城。汉王至河南。韩信急击韩王上阳城。上降。汉王乃立韩信为韩王。常将韩兵从。三年。汉王出荥阳。韩王信。周苛等守荥阳。及楚败荥阳。信降楚。已而得亡。复归汉。汉复立以为韩王。竟从击破项籍。天下定。五年春。遂与剖符为韩王。王颍川。\\
\hspace*{2em}明年春。上以韩信材武。所王北近巩。洛。南迫宛。叶。东有淮阳。皆天下劲兵处。乃诏徙韩王信王太原以北。备御胡。都晋阳。信上书曰。﹁国被边。匈护数入。晋阳去塞远。请治马邑。﹂上许之。信乃徙治马邑。秋。匈护冒顿大围信。信数使使胡求和解。汉发兵救之。疑信数间使。有二心。使人责让信。信恐诛。因与匈护约共攻汉。反。以马邑降胡。击太原。\\
\hspace*{2em}七年冬。上自往击。破信军铜鞮。斩其将王喜。信亡走匈护。其与白土人曼丘臣。王黄等立赵苗裔赵利为王。复收信败散兵。而与信及冒顿谋攻汉。匈护仗左右贤王将万馀骑与王黄等屯广武以南。至晋阳。与汉兵战。汉大破之。追至于离石。破之。匈护复聚兵楼烦西北。汉令车骑击破匈护。匈护常败走。汉乘胜追北。闻冒顿居代谷。高皇帝居晋阳。使人视冒顿。还报曰﹁可击﹂。上遂至平城。上出白登。匈护骑围上。上乃使人厚遗阏氏。阏氏乃说冒顿曰。﹁今得汉地。犹不能居。且两主不相戹。﹂居七日。胡骑稍将去。时天大雾。汉使人往来。胡不觉。护军中尉陈平言上曰。﹁胡者全兵。请令彊弩傅两矢外乡。徐行出围。﹂入平城。汉救兵亦到。胡骑遂解去。汉亦罢兵归。韩信为匈护将兵往来击边。\\
\hspace*{2em}汉十年。信令王黄等说误陈豨。十一年春。故韩王信复与胡骑入居参合。距汉。汉使柴将军击之。遗信书曰。﹁陛下宽仁。诸侯虽有畔亡。而复归。辄复故位号。不诛也。大王所知。今王以败亡走胡。非有大罪。急自归。﹂韩王信报曰。﹁陛下擢仆起闾巷。南面称孤。此仆之幸也。荥阳之事。仆不能死。囚於项籍。此一罪也。及寇攻马邑。仆不能坚守。以城降之。此二罪也。今反为寇将兵。与将军争一旦之命。此三罪也。夫种。蠡成一罪。身死亡。今仆有三罪於陛下。而欲求活於世。此伍子胥所以偾於吴也。今仆亡匿山谷间。旦暮乞贷蛮夷。仆之思归。如痿人不忘起。盲者不忘视也。势不可耳。﹂遂战。柴将军屠参合。斩韩王信。\\
\hspace*{2em}信之入匈护。与太子俱。及至穨当城。生子。因名曰穨当。韩太子亦生子。命曰婴。至孝文十四年。穨当及婴率其众降汉。汉封穨当为弓高侯。婴为襄城侯。吴楚军时。弓高侯功冠诸将。传子至孙。孙成子。失侯。婴孙以不敬失侯。穨当孽孙韩嫣。贵幸。名富显於当世。其弟说。再封。数称将军。卒为案道侯。子代。岁馀坐法死。後岁馀。说孙曾拜为龙嵒侯。续说後。\\
\hspace*{2em}卢绾者。丰人也。与高祖同里。卢绾亲与高祖太上皇相爱。及生男。高祖。卢绾同日生。里中持羊酒贺两家。及高祖。卢绾壮。俱学书。又相爱也。里中嘉两家亲相爱。生子同日。壮又相爱。复贺两家羊酒。高祖为布衣时。有吏事辟匿。卢绾常随出入上下。及高祖初起沛。卢绾以客从。入汉中为将军。常侍中。从东击项籍。以太尉常从。出入卧内。衣被饮食赏赐。群臣莫敢望。虽萧曹等。特以事见礼。至其亲幸。莫及卢绾。绾封为长安侯。长安。故咸阳也。\\
\hspace*{2em}汉五年冬。以破项籍。乃使卢绾别将。与刘贾击临江王共尉。破之。七月还。从击燕王臧荼。臧荼降。高祖已定天下。诸侯非刘氏而王者七人。欲王卢绾。为群臣觖望。及虏臧荼。乃下诏诸将相列侯。择群臣有功者以为燕王。群臣知上欲王卢绾。皆言曰。﹁太尉长安侯卢绾常从平定天下。功最多。可王燕。﹂诏许之。汉五年王月。乃立卢绾为燕王。诸侯王得幸莫如燕王。\\
\hspace*{2em}汉十一年秋。陈豨反代地。高祖如邯郸击豨兵。燕王绾亦击其东北。当是时。陈豨使王黄求救匈护。燕王绾亦使其臣张胜於匈护。言豨等军破。张胜至胡。故燕王臧茶子衍出亡在胡。见张胜曰。﹁公所以重於燕者。以习胡事也。燕所以久存者。以诸侯数反。兵连不决也。今公为燕欲急灭豨等。豨等已尽。次亦至燕。公等亦且为虏矣。公何不令燕且缓陈豨而与胡和。事宽。得长王燕。即有汉急。可以安国。﹂张胜以为然。乃私令匈护助豨等击燕。燕王绾疑张胜与胡反。上书请族张胜。胜还。具道所以为者。燕王寤。乃诈论它人。脱胜家属。使得为匈护间。而阴使范齐之陈豨所。欲令久亡。连兵勿决。\\
\hspace*{2em}汉十二年。东击黥布。豨常将兵居代。汉使樊哙击斩豨。其裨将降。言燕王绾使范齐通计谋於豨所。高祖使使召卢绾。绾称病。上又使辟阳侯审食其。御史大夫赵尧往迎燕王。因验问左右。绾愈恐。闭匿。谓其幸臣曰。﹁非刘氏而王。独我与长沙耳。往年春。汉族淮阴。夏。诛彭越。皆吕后计。今上病。属任吕后。吕后妇人。专欲以事诛异姓王者及大功臣。﹂乃遂称病不行。其左右皆亡匿。语颇泄。辟阳侯闻之。归具报上。上益怒。又得匈护降者。降者言张胜亡在匈护。为燕使。於是上曰。﹁卢绾果反矣。﹂使樊哙击燕。燕王绾悉将其宫人家属骑数千居长城下。侯伺。幸上病愈。自入谢。四月。高祖崩。卢绾遂将其众亡入匈护。匈护以为东胡卢王。绾为蛮夷所侵夺。常思复归。居岁馀。死胡中。\\
\hspace*{2em}高后时。卢绾妻子亡降汉。会高后病。不能见。舍燕邸。为欲置酒见之。高后竟崩。不得见。卢绾妻亦病死。\\
\hspace*{2em}孝景中六年。卢绾孙他之。以东胡王降。封为亚谷侯。\\
\hspace*{2em}陈豨者。宛朐人也。不知始所以得从。及高祖七年冬。韩王信反。入匈护。上至平城还。乃封豨为列侯。以赵相国将监赵。代边兵。边兵皆属焉。\\
\hspace*{2em}豨常告归过赵。赵相周上见豨宾客随之者千馀乘。邯郸官舍皆满。豨所以待宾客布衣交。皆出客下。豨还之代。周上乃求入见。见上。具言豨宾客盛甚。擅兵於外数岁。恐有变。上乃令人覆案豨客居代者财物诸不法事。多连将豨。豨恐。阴令客通使王黄。曼丘臣所。及高祖十年七月。太上皇崩。使人召豨。豨称病甚。九月。遂与王黄等反。自立为代王。劫略赵。代。\\
\hspace*{2em}上闻。乃赦赵。代吏人为豨所诖误劫略者。皆赦之。上自往。至邯郸。喜曰。﹁豨不南据漳水。北守邯郸。知其成能为也。﹂赵相奏斩常山守。尉。曰。﹁常山二十五城。豨反。亡其二十城。﹂上问曰。﹁守。尉反乎。﹂对曰。﹁不反。﹂上曰。﹁是力不足也。﹂赦之。复以为常山守。尉。上问周上曰。﹁赵亦有壮士可令将者乎。﹂对曰。﹁有四人。﹂四人谒。上谩骂曰。﹁竖子能为将乎。﹂四人惭伏。上封之各千户。以为将。左右谏曰。﹁从入蜀。汉。伐楚。功未遍行。今此何功而封。﹂上曰。﹁非若所知。陈豨反。邯郸以北皆豨有。吾以羽檄徵天下兵。未有至者。今唯独邯郸中兵耳。吾胡爱四千户封四人。不以慰赵子弟。﹂皆曰。﹁善。﹂於是上曰。﹁陈豨将谁。﹂曰。﹁王黄。曼丘臣。皆故贾人。﹂上曰。﹁吾知之矣。﹂乃各以千金购黄。臣等。\\
\hspace*{2em}十一年冬。汉兵击斩陈豨将侯敞。王黄於曲逆下。破豨将张春於聊城。斩首万馀。太尉勃入定太原。代地。十二月。上自击东垣。东垣不下。卒骂上。东垣降。卒骂者斩之。不骂者黥之。更命东垣为真定。王黄。曼丘臣其麾下受购赏之。皆生得。以故陈豨军遂败。\\
\hspace*{2em}上还至洛阳。上曰。﹁代居常山北。赵乃从山南有之。远。﹂乃立子恆为代王。都中都。代。雁门皆属代。\\
\hspace*{2em}高祖十二年冬。樊哙军卒追斩豨於灵丘。\\
\hspace*{2em}太史公曰。韩信。卢绾非素积德累善之世。徼一时权变。以诈力成功。遭汉初定。故得列地。南面称孤。内见疑彊大。外倚蛮貊以为援。是以日疏自危。事穷智困。卒赴匈护。岂不哀哉。陈豨。梁人。其少时数称慕魏公子。及将军守边。招致宾客而下士。名声过实。周上疑之。疵瑕颇起。惧祸及身。邪人进说。遂陷成道。於戏悲夫。夫计之生孰成败於人也深矣。\\

% ===== 索隐述赞（司马贞《史记索隐》）· 统一样式 =====
% 题头「索隐述赞曰」五字并入反白框；正文每四句一列，定宽盒 \hfil 居列中；
% \fontsize{24}{41} 列距对齐全书界行（121+41k）；\clearpage 使述赞独占一页。
\clearpage
\shuzanhead{索隐述赞曰}\linebreak
{\shuzanbody\fontsize{24}{41}\selectfont
\hbox to \hsize{\hfil 韩襄遗孽。始从汉中。剖符南面。徙邑北通。\hfil}\linebreak
\hbox to \hsize{\hfil 穨当归国。龙雒有功。卢绾亲爱。群臣莫同。\hfil}\linebreak
\hbox to \hsize{\hfil 旧燕是王。东胡计穷。\hfil}%
\par}
\end{document}
